With the provisioning services enabled, the next step is to define and
customize a system image that can subsequently be used to provision one or more
{\em compute} nodes. The following subsections highlight this process.

\subsubsection{Build initial BOS image} \label{sec:assemble_bos}
\Warewulf{} 4 supports using container images as the base file system for
provisioning, and it can import these directly from an OCI registry like
Docker Hub. In this example we import the official Ubuntu container image and
install the kernel and init-system packages needed for Warewulf provisioning.

The \texttt{wwctl image exec} command runs the commands below it, these commands
also be run interactively one a time with the command \texttt{wwctl image
shell \baseos{}}. You can add \texttt{/bin/false} as the last command to prevent
the image from rebuilding (it will show an error) and rebuild later with the
\texttt{wwctl image build} command.

% begin_ohpc_run
% ohpc_comment_header Create compute image for Warewulf \ref{sec:assemble_bos}
\begin{lstlisting}[language=bash,literate={-}{-}1,keywords={},upquote=true,keepspaces,literate={BOSVER}{\baseos{}}1 {OSIMAGE}{\osimage{}}1]
# Import the base image
[sms](*\#*) wwctl image import docker://OSIMAGE BOSVER --syncuser

# Define chroot location
[sms](*\#*) export CHROOT=$(wwctl image show BOSVER)

# Enable OpenHPC inside image and update image. Disable image build on exit, rebuild later.
[sms](*\#*) wwctl image exec --build=false BOSVER -- /bin/bash -ex <<- EOF
  export DEBIAN_FRONTEND=noninteractive
  apt update
  apt -y install ca-certificates curl gnupg systemd systemd-sysv \
    udev kmod linux-image-generic initramfs-tools network-manager openssh-server
  OHPC_REPO_ROOT=https://repos.versatushpc.com.br/openhpc/versatushpc-4
  OHPC_UBUNTU_REPO=${OHPC_REPO_ROOT}/Ubuntu_24.04
  install -d -m 0755 /usr/share/keyrings /etc/apt/sources.list.d
  curl -fsSL ${OHPC_REPO_ROOT}/versatushpc.gpg \
    -o /usr/share/keyrings/versatushpc.gpg
  curl -fsSL ${OHPC_UBUNTU_REPO}/versatushpc-openhpc.list \
    -o /etc/apt/sources.list.d/versatushpc-openhpc.list
  apt update
  apt -y upgrade
EOF
\end{lstlisting}
% end_ohpc_run
